\section{Метод}

\subsection{Базовая модель следующего айтема}

Политика --- декодер только с самовниманием~\cite{gpt2}: 4 слоя, 4 головы,
размерность 128, контекст 200 событий, связанные эмбеддинги входа и выхода.
При словаре Listen+ из 630\,974 айтемов модель насчитывает 81.6 млн
параметров, из которых 80.8 млн приходятся на таблицу айтемов. История
строится по всем прослушиваниям подряд, включая отклонённые: контекст модели
не разрывается на пропусках, поэтому переходы между треками остаются
реалистичными.

Обучающий сигнал --- следующий айтем для каждой позиции окна, как в
\code{TrainDataset} SASRec~\cite{sasrec}, а не один target на окно. Для
случайного пользователя выбирается случайный конец окна, из истории берутся
последние 200 событий, и таргетами становятся все следующие позиции окна.
Полный softmax по 630\,974 айтемам для батча 32 и 200 позиций занял бы
около 16 ГБ в fp16 только на логиты, поэтому обучение идёт с
\emph{sampled softmax}:
к целевым айтемам добавляются 8192 случайных негатива из униформного
распределения по каталогу, и минимизируется
\begin{equation}
\mathcal{L} = -\log \frac{e^{s_y}}{\,e^{s_y} + \sum_{j\in\mathcal{N}} e^{s_j}\,},
\end{equation}
где $s$ --- логит айтема. In-batch негативы берутся из уникальных целей
окна (до 6400 на батч). Оценка при этом всегда идёт по полному каталогу,
так что метрики не завышаются выбором кандидатов.

Оптимизация: AdamW, $\text{lr}=3\cdot10^{-4}$, прогрев 100 шагов,
косинусное затухание до $3\cdot10^{-5}$, weight decay $0.01$, обрезка нормы
градиента 1.0, fp16, 3000 шагов (около 4 минут на T4; при 32 окнах по 200
позиций это порядка 14 млн обучающих пар). Чекпоинт выбирается по NDCG@10
на валидационном дне каждые 500 шагов.

\subsection{Reward model}
\label{sec:method-rm}

Reward model оценивает пару (история, кандидат). Вход --- последние 100
событий истории Listen+, представленных эмбеддингом айтема, плюс вектор
признаков кандидата (логарифм популярности, средняя длина трека, доля
лайков) и сводные признаки истории (средняя доля прослушивания, доля
органических событий, длина). Кандидат вставляется последним токеном
последовательности, его представление и есть вход голов. Кодировщик ---
2 слоя, 4 головы, размерность 128; маска закрывает паддинги истории, но не
кандидата.

Четыре головы:
\begin{itemize}[leftmargin=1.2em,itemsep=1pt,topsep=2pt]
  \item $r_{\text{ratio}}$ --- доля досмотра. Распределение этой величины
  у Listen+ сосредоточено в $[0.5,1]$, у отклонённых прослушиваний ---
  в $[0,0.5)$. Чтобы модель умела штрафовать отклонённый трек, голова
  несёт два выхода: бинарный «прослушано не меньше половины» и
  регрессию через сигмоиду на нормированной доле досмотра. Цель регрессии
  клипуется в $[0,1]$, так как 0.47\% событий имеют
  \code{played\_ratio\_pct} $>100$ при повторном прослушивании;
  \item $r_{\text{like}}$ --- лайк на этот айтем позже внутри сплита
  (\code{like\_next}). Состояние лайка --- накопленная сумма like/unlike,
  поэтому поздний unlike отменяет лайк;
  \item $r_{\text{sess}}$ --- продолжение сессии: следующее событие
  пользователя не позже 30 минут;
  \item $r_{\text{ret}}$ --- возврат пользователя в течение 7 дней.
\end{itemize}

Выборка обучающего периода несбалансирована по замыслу: 40\% реальные
события Listen+, 25\% hard negatives (другой айтем из истории того же
пользователя, который он \emph{не} дослушал до половины), 20\% скипы
(события с $0 <$ доля досмотра $< 0.5$, с настоящей целью регрессии) и
15\% случайные айтемы каталога с ненулевой популярностью. На одну модель
приходится 1.5 млн пар, один проход по 256 пар на шаг. Общая потеря:
\begin{equation}
\mathcal{L} =
\mathrm{BCE}_{\text{listen}}
+\mathrm{MSE}_{[0,1]}
+0.5\,\mathrm{BCE}_{\text{sess}}
+0.5\,\mathrm{BCE}_{\text{ret}}
+\mathrm{BCE}_{\text{like}},
\end{equation}
где $\mathrm{BCE}_{\text{listen}}$ --- кросс-энтропия бинарного выхода
головы досмотра на всех ветках, а $\mathrm{MSE}_{[0,1]}$ --- её регрессия
только на реальных событиях и скипах, у которых есть наблюдённая доля
досмотра. Веса вспомогательных голов снижены: они
зависят от пользователя и места в истории, а не от кандидата, и при весе 0.5
доминировали в общей потере.

Ансамбль из трёх моделей обучается с разными seed на подвыборках
пользователей (80\% пользователей на модель), поэтому разброс оценок между
моделями имеет поведенческий смысл: он отражает разницу в исходных
подвыборках, а не шум инициализации. Награда в RL --- пессимистичная оценка
\begin{equation}
R = \bar r - \beta\,\sigma_r,\qquad \beta = 0.5,
\end{equation}
где $\bar r$ и $\sigma_r$ --- среднее и стандартное отклонение по ансамблю;
это стандартный приём против переоптимизации суррогата~\cite{gao2022}.
Пессимизм применяется к каждой голове отдельно; в blend-награде
$\bar r = \bar r_{\text{listen}} + \lambda\,\bar r_{\text{like}}$.

\subsection{Офлайн-обучение с подкреплением}

Основной метод --- GRPO~\cite{grpo} в клиппованной форме. Для одного промпта
(истории дня валидации) политика порождает группу из $G=8$ роллаутов по
$H=5$ айтемов; возврат роллаута $R_i$ --- сумма наград с дисконтом
$\gamma=1$. Преимущество нормируется внутри группы,
\begin{equation}
A_i = \frac{R_i - \operatorname{mean}(R_{1:G})}{\operatorname{std}(R_{1:G}) + \varepsilon},
\qquad |A_i| \le 5.
\end{equation}
Каждый шаг оптимизации использует отношение
$\rho_i = \pi_\theta(a_i)/\pi_{\theta_{\text{old}}}(a_i)$ и клиппованный
суррогат
\begin{equation}
\mathcal{J} = -\frac{1}{G}\sum_{i=1}^{G}
\min\!\left(\rho_i A_i,\ \operatorname{clip}(\rho_i,1-\epsilon,1+\epsilon)A_i\right)
\;+\; \beta_{\text{KL}}\,\widehat{\mathrm{KL}}\!\left(\pi_\theta \,\|\, \pi_{\text{ref}}\right),
\end{equation}
$\epsilon=0.2$, один проход на батч, $\beta_{\text{KL}}=0.05$. KL оценивается
k3-формулой $e^{-d} + d - 1$ при $d = \log\pi_\theta - \log\pi_{\text{ref}}$:
при сэмплах из $\pi_\theta$ это несмещённая оценка
$\mathrm{KL}(\pi_\theta\|\pi_{\text{ref}})$, неотрицательная поэлементно.
Параметр $\theta$ инициализируется базовой моделью, $\pi_{\text{ref}}$
заморожена. Айтемы сэмплируются из полного softmax по каталогу, а не из
top-100, поэтому поведенческая и обучаемая политики совпадают.

Кроме GRPO, в коде реализованы DPO~\cite{dpo} на парах роллаутов из RM и
filtered behavior cloning по роллаутам с наградой выше медианы; в этой работе
основной метод --- GRPO с наградой по голове досмотра.

\paragraph{Среда.} Симулятор разомкнутый: пять сгенерированных айтемов не
меняют будущую историю и не влияют на последующие награды. Это главное
упрощение метода: долгосрочные эффекты (H5) такой симулятор не различает,
а оптимизация работает как одношаговая переоценка пар (история, айтем).
